\chapter{\ifenglish Introduction\else บทนำ\fi}

\section{\ifenglish Project rationale\else ที่มาของโครงงาน\fi}
\enskip ปัจจุบันมีการส่งเสริมให้นักศึกษามหาวิทยาลัยเชียงใหม่หันมาใส่ใจสุขภาพและการออกกำลังกายมากขึ้น \\โดยมีบริการฟิตเนสภายในมหาวิทยาลัย เช่น CMU Fitness by Playground อย่างไรก็ตาม พบว่าผู้ใช้งานจำนวนไม่น้อยยังประสบปัญหาขาดแรงจูงใจ ความต่อเนื่องในการออกกำลังกาย รวมถึงความยุ่งยากในการวางแผนการฝึกและการเลือกใช้อุปกรณ์อย่างเหมาะสม

\enskip จากสภาพปัญหาดังกล่าว จึงเกิดแนวคิดในการพัฒนาโครงงาน GymMate ซึ่งมีจุดมุ่งหมายเพื่อเป็นต้นแบบระบบที่ช่วยสนับสนุนการออกกำลังกายของนักศึกษา โดยเน้นการสร้างแรงจูงใจและการจัดการการใช้งานฟิตเนสให้มีประสิทธิภาพมากยิ่งขึ้น

\section{\ifenglish Objectives\else วัตถุประสงค์ของโครงงาน\fi}
\setlength{\leftmargini}{1.25cm}
\begin{enumerate}
\renewcommand{\labelenumi}{1.2.\arabic{enumi}}
    \item พัฒนาระบบการวางแผนการออกกำลังกาย (Workout Program Planner)\\
    GymMate มุ่งเน้นการพัฒนาโปรแกรมฝึกที่สามารถเลือกได้ทั้งรูปแบบสำเร็จรูป (preset) และการปรับแต่งตามความต้องการ (custom) เพื่อช่วยให้นักศึกษาสามารถวางแผนการออกกำลังกาย\\ได้สะดวกและเหมาะสมกับเป้าหมายของตนเอง

    \item เพิ่มประสิทธิภาพการออกกำลังกายด้วยระบบ Flex Substitute\\
    ด้วยระบบแนะนำท่าและเครื่องเล่นทดแทนอัตโนมัติ เมื่ออุปกรณ์ที่ต้องการไม่พร้อมใช้งาน ผู้ใช้สามารถเลือกทางเลือกที่เหมาะสมกับกล้ามเนื้อที่ต้องการฝึกได้ โดยระบบจะช่วยแนะนำผ่านอัลกอริทึมที่ออกแบบมาโดยเฉพาะ

    \item จัดการฟิตเนสอย่างเป็นระบบผ่าน CMS Dashboard\\
    ฝั่งผู้ดูแลฟิตเนสสามารถจัดการตำแหน่งและข้อมูลเครื่องเล่นได้อย่างมีประสิทธิภาพ \\ลดความซับซ้อนในการบริหารจัดการ และเพิ่มความสะดวกให้กับทั้งผู้ดูแลและผู้ใช้งาน

    \item สร้างแรงจูงใจและความต่อเนื่องในการออกกำลังกาย\\
    ด้วยระบบ Gamification และ Leaderboard ผู้ใช้งานสามารถสะสมคะแนน แข่งขันกับผู้ใช้คนอื่น ติดตามความก้าวหน้าของตนเอง ทำให้การออกกำลังกายน่าสนใจ สนุก และมีแรงกระตุ้นมากขึ้น
\end{enumerate}

\section{\ifenglish Project scope\else ขอบเขตของโครงงาน\fi}
\enskip โครงงาน GymMate มีขอบเขตการดำเนินงานที่เน้นการ สำรวจและออกแบบระบบต้นแบบ โดยมีรายละเอียดดังนี้

\subsection{\ifenglish  Target Users\else กลุ่มเป้าหมาย\fi}
\begin{itemize}
    \item \textbf{ผู้ใช้บริการฟิตเนส (Fitness Members):}\\
    โดยเฉพาะกลุ่มผู้เริ่มต้น (Beginners) และนักศึกษามหาวิทยาลัยเชียงใหม่ ที่ต้องการโปรแกรมฝึก การติดตามความก้าวหน้า และแรงจูงใจในการออกกำลังกาย
    \item \textbf{ผู้จัดการและเจ้าของฟิตเนส (Gym Managers \&  Owners):}\\
    ต้องการเครื่องมือในการบริหารจัดการเครื่องเล่นและสื่อสารกับสมาชิก เพื่อเพิ่มประสิทธิภาพการใช้งานพื้นที่และเสริมความพึงพอใจของลูกค้า
\end{itemize}

\subsection{\ifenglish  System Features in Scope\else ขอบเขตฟังก์ชัน\fi}
\begin{enumerate}
    \item \textbf{Workout Program Planner} ระบบวางแผนโปรแกรมการออกกำลังกาย\\
    ผู้เริ่มต้นสามารถเลือกใช้ Preset Plan เพื่อสร้างพื้นฐานการฝึกอย่างเป็นระบบ หรือปรับเป็น Custom Plan ให้สอดคล้องกับเป้าหมายเฉพาะบุคคล พร้อมแนวทาง Goal \& Approach เช่น Push–Pull–Legs, Upper–Lower หรือ Total Body เพื่อเพิ่มประสิทธิภาพ
    \item \textbf{Floorplan UI \& Flex Substitute} ส่วนติดต่อผู้ใช้สำหรับผังพื้นที่และการปรับเปลี่ยนแบบยืดหยุ่น\\
    ระบบมีการจัดทำ แผนผังยิม เพื่อแสดงตำแหน่งและการจัดวางของเครื่องเล่นอย่างชัดเจน อีกทั้งมีฟังก์ชัน ปุ่ม Flex Substitute ที่ช่วยแนะนำเครื่องหรือท่าทางทดแทนเมื่ออุปกรณ์ที่ต้องการใช้งานเต็ม โดยยังคงมุ่งเน้นการฝึกที่กล้ามเนื้อเป้าหมายเดิมเพื่อไม่ให้ผลลัพธ์การออกกำลังกายคลาดเคลื่อน.
    \item \textbf{User Management \& Tracking} ระบบจัดการและติดตามผู้ใช้งาน\\
    ระบบมี โปรไฟล์ผู้ใช้ สำหรับจัดเก็บข้อมูลการฝึก พร้อมการบันทึกสถิติ เช่น จำนวนเซต, ครั้งทำซ้ำ และน้ำหนักที่ใช้ นอกจากนี้ยังมี กราฟความก้าวหน้า เพื่อสะท้อนพัฒนาการอย่างต่อเนื่อง รวมถึงรายงานผลการฝึกในรูปแบบรายสัปดาห์และรายเดือน เพื่อช่วยประเมินและปรับแผนการออกกำลังกายได้อย่างเป็นระบบ
    \item \textbf{Gamification Layer} ชั้นการใช้องค์ประกอบเชิงเกม\\
    ระบบมีการสะสมคะแนนจากความสม่ำเสมอในการฝึก พร้อมกลไกสร้างแรงจูงใจผ่าน Badge, Level, Streak และ Leaderboard ซึ่งช่วยเสริมทั้งความต่อเนื่องและการแข่งขันเชิงบวกระหว่างผู้ใช้
    \item \textbf{Content \& Notification} ระบบจัดการเนื้อหาและการแจ้งเตือน\\
    ระบบรองรับ วิดีโอและภาพประกอบท่าทาง ผ่าน QR code ที่ติดไว้บนเครื่องออกกำลังกาย เพื่อให้ผู้ใช้เข้าถึงวิธีการใช้งานได้อย่างถูกต้อง นอกจากนี้ยังมี การแจ้งเตือน (Notification) ตามตารางโปรแกรม รวมถึงการเปลี่ยนท่าอัตโนมัติเมื่อเครื่องที่ต้องการใช้งานเต็ม เพื่อคงประสิทธิภาพของการฝึก
    \item \textbf{Admin Dashboard (CMS)} แดชบอร์ดผู้ดูแลระบบ\\
    ระบบสามารถ จัดการข้อมูลเครื่องเล่น ได้อย่างยืดหยุ่น ทั้งการเพิ่ม ย้าย หรือแก้ไขสถานะของเครื่องออกกำลังกาย พร้อมทั้งแสดง ข้อมูลพื้นฐานของฟิตเนส เช่น เวลาเปิด–ปิด และข่าวสารที่เกี่ยวข้อง เพื่ออำนวยความสะดวกแก่ผู้ใช้
\end{enumerate}

\subsection{\ifenglish Hardware scope\else ขอบเขตด้านฮาร์ดแวร์\fi}
\enskip โครงงานนี้ไม่ได้พัฒนาฮาร์ดแวร์ใหม่ แต่ใช้โครงสร้างพื้นฐานและบริการที่มีอยู่ ได้แก่
\begin{itemize}
    \item สมาร์ทโฟนของผู้ใช้งานสำหรับติดตั้งและใช้งานแอปพลิเคชัน
    \item สถานที่และเครื่องเล่นออกกำลังกายของยิม
\end{itemize}

\subsection{\ifenglish Software scope\else ขอบเขตด้านซอฟต์แวร์\fi}
\enskip ขอบเขตด้านซอฟต์แวร์ของโครงงานประกอบด้วย
\begin{itemize}
    \item ต้องรองรับการใช้งานและติดตั้งแอปพลิเคชันบนสมาร์ทโฟนของผู้ใช้
    \item ต้องรองรับการใช้งานผ่านคอมพิวเตอร์สำหรับการจัดการ CMS Dashboard ของผู้ดูแลระบบ
\end{itemize}


\section{\ifenglish Expected outcomes\else ประโยชน์ที่ได้รับ\fi}
\setlength{\leftmargini}{1.25cm}
\begin{enumerate}
\renewcommand{\labelenumi}{1.4.\arabic{enumi}}
    \item ได้ต้นแบบแอปพลิเคชัน GymMate ที่ช่วยสนับสนุนการออกกำลังกายและเพิ่มแรงจูงใจสำหรับนักศึกษามหาวิทยาลัยเชียงใหม่
    \item ผู้ใช้งานสามารถวางแผนโปรแกรมการออกกำลังกายได้สะดวกและเหมาะสมกับเป้าหมายของตนเอง ทั้งแบบสำเร็จรูปและแบบปรับแต่ง
    \item ผู้ดูแลฟิตเนสสามารถบริหารจัดการตำแหน่งและข้อมูลเครื่องเล่นได้ง่ายขึ้น ผ่าน CMS Dashboard
    \item ระบบช่วยเพิ่มความต่อเนื่อง สนุก และประสิทธิภาพในการออกกำลังกายของนักศึกษา
    \item โครงงานสามารถต่อยอดไปสู่การใช้งานจริง และขยายผลสำหรับฟิตเนสหรือมหาวิทยาลัยอื่นได้
\end{enumerate}

\section{\ifenglish Technology and tools\else เทคโนโลยีและเครื่องมือที่ใช้\fi}
\enskip เพื่อพัฒนาและนำเสนอระบบ GymMate ได้ตามเป้าหมาย โครงงานนี้เลือกใช้เทคโนโลยีหลัก 3 ส่วน ดังนี้

\begin{enumerate}
\renewcommand{\labelenumi}{1.5.\arabic{enumi}}
    \item \textbf{Web Application (CMS Dashboard)}\\
    ระบบสำหรับ ผู้ดูแล (Admin/Gym Manager) ทำหน้าที่บริหารจัดการข้อมูลเครื่องเล่น รวมถึงการย้ายหรือแก้ไขสถานะใน Floorplan และจัดการข้อมูลฟิตเนส เช่น ข่าวสารและเวลาเปิด–ปิด เพื่อให้การดำเนินงานเป็นระเบียบและมีประสิทธิภาพ
    \item \textbf{Mobile Application (User App)}\\
    ระบบสำหรับ ผู้ใช้ฟิตเนส (Members) รองรับการเลือกโปรแกรมฝึกที่เหมาะสม \\พร้อมฟังก์ชัน Calendar/Day Plan เพื่อดูตารางการฝึก และ Floorplan Highlight ช่วยระบุตำแหน่งเครื่องเล่น นอกจากนี้ยังมี Flex Substitute สำหรับแนะนำท่าทางทดแทน, Log Progress เพื่อบันทึกผลการฝึก, ระบบ Gamification เพื่อสร้างแรงจูงใจ และ Notification สำหรับแจ้งเตือนตามโปรแกรมอย่างต่อเนื่อง
    \item \textbf{Server \& Backend Services}\\
    ระบบนี้ใช้สำหรับ ประมวลผลและจัดเก็บข้อมูล ที่เกี่ยวข้องกับโปรแกรมการฝึก การใช้งาน และการจัดการเครื่อง โดยมีฟังก์ชันหลัก ได้แก่ Workout Planner สำหรับวางแผนการฝึก, Flex Substitute Algorithm เพื่อแนะนำทางเลือกเมื่อเครื่องเต็ม, ระบบ Authentication/Authorization เพื่อควบคุมสิทธิ์การใช้งาน และ Data Logging สำหรับเก็บบันทึกข้อมูลอย่างเป็นระบบ
\end{enumerate}

\section{\ifenglish Project plan\else แผนการดำเนินงาน\fi}

\begin{plan}{6}{2025}{3}{2026}
    \planitem{6}{2025}{7}{2025}{ศึกษาค้นคว้าและหาหัวข้อสําหรับโครงงาน}
    \planitem{7}{2025}{9}{2025}{เก็บและวิเคราะห์ความต้องการ}
    \planitem{8}{2025}{9}{2025}{ออกแบบฟีเจอร์และปรับขอบเขต}
    \planitem{9}{2025}{11}{2025}{ออกแบบ UX/UI และต้นแบบ}
    \planitem{11}{2025}{12}{2025}{ออกแบบเชิงเทคนิค: Database, API, System Architecture}
    \planitem{12}{2025}{2}{2026}{พัฒนา Backend}
    \planitem{1}{2026}{2}{2026}{พัฒนา Mobile App}
    \planitem{1}{2026}{2}{2026}{พัฒนา Web CMS (Admin Dashboard, Floorplan Editor)}
    \planitem{2}{2026}{3}{2026}{ทดลองใช้งานจริงกับผู้ใช้}
\end{plan}

\section{\ifenglish Roles and responsibilities\else บทบาทและความรับผิดชอบ\fi}
\begin{enumerate}
    \item นายภาณุเดช เสือเผือก รับบทบาท Product Owner, Frontend
    \begin{itemize}
        \item \textbf{หน้าที่:} เสนอไอเดียจากประสบการณ์การเข้ายิม, พัฒนา Mobile/Web Frontend, ตรวจสอบความถูกต้องของข้อมูลการออกกำลังกาย
        \item \textbf{ความรู้ที่ใช้:} ความรู้ด้านการออกกำลังกาย และการพัฒนาส่วนติดต่อผู้ใช้ (Frontend Development)
    \end{itemize}

    \item นายภูเบศร์ เรืองคุ้ม รับบทบาท Backend \& System
    \begin{itemize}
        \item \textbf{หน้าที่:} ออกแบบสถาปัตยกรรมระบบและฐานข้อมูล, พัฒนา API และ Flex Substitute และความปลอดภัยระบบ
        \item \textbf{ความรู้ที่ใช้:} การเขียนโปรแกรมฝั่งเซิร์ฟเวอร์ (Backend Programming), ฐานข้อมูล และความปลอดภัยของระบบ
    \end{itemize}
    
    \item นายอนวัช รัตนกิจศร รับบทบาท UX/UI \& Documentation
    \begin{itemize}
        \item \textbf{หน้าที่:} ออกแบบ UX/UI Mockup, จัดทำเอกสารรายงาน, ประสานงานกับอาจารย์และ Stakeholder, สนับสนุนงาน Frontend
        \item \textbf{ความรู้ที่ใช้:} การออกแบบประสบการณ์ผู้ใช้และส่วนติดต่อผู้ใช้ (UX/UI Design), การเขียนเชิงเทคนิค (Technical Writing), การสื่อสาร (Communication)
    \end{itemize}
    
\end{enumerate}

\section{\ifenglish%
Impacts of this project on society, health, safety, legal, and cultural issues
\else%
ผลกระทบด้านสังคม สุขภาพ ความปลอดภัย กฎหมาย และวัฒนธรรม
\fi}

\enskip โครงงาน GymMate อาจมีผลกระทบและการประยุกต์ใช้งานในหลายมิติ ดังนี้

\begin{enumerate}
    \item \textbf{ด้านสังคม:} สนับสนุนให้นักศึกษาหันมาออกกำลังกายมากขึ้น 
    ช่วยสร้างแรงจูงใจผ่านระบบ Gamification และ Leaderboard
    \item \textbf{ด้านสุขภาพ:} ช่วยส่งเสริมสุขภาพกายและจิตใจของผู้ใช้ 
    ลดความเสี่ยงจากโรคที่เกิดจากการขาดการออกกำลังกาย
    \item \textbf{ด้านความปลอดภัย:} การแนะนำการใช้อุปกรณ์และท่าทางอย่างถูกต้อง 
    ลดความเสี่ยงในการบาดเจ็บจากการออกกำลังกาย
    \item \textbf{ด้านกฎหมาย:} ต้องคำนึงถึงลิขสิทธิ์ของข้อมูล รูปภาพ และเพลง 
    รวมถึงการจัดเก็บข้อมูลผู้ใช้ให้เป็นไปตามกฎหมายคุ้มครองข้อมูลส่วนบุคคล (PDPA/GDPR)
    \item \textbf{ด้านวัฒนธรรม:} ปรับการออกแบบแอปให้เหมาะกับวิถีชีวิตและพฤติกรรมของนักศึกษา 
    ส่งเสริมการสร้างวัฒนธรรมการดูแลสุขภาพในมหาวิทยาลัย
\end{enumerate}
